% \chapter{Лекция}
% \label{ch:intro}

% \section{Постановка проблемы}
% После формирования передаваемых символов в программе на ПК отсчеты символов передаются на SDR, где происходит формирование радиосигнала
% на заданной частоте несущего колебания. На приемной стороне происходит обратное преобразование с радиочастоты до получения цифровых отсчетов. При
% этом частота генератора несущего колебания передатчика $f_{ct}$ и частота генератора несущего опорного колебания приемника $f_{cr}$ хоть и задаются в программе
% одинаково, но фактически в SDR они могут отличаться, отличие в значении
% несущего колебания может достигать нескольких кГц. Это отличие называется
% смещением несущего колебания.

% $$f_o = f_{ct} - f_{cr}$$

% Частотное смещение проявляется в повороте принятого комплексного НЧ
% сигнала на величину $e^{i2\pi f_ot}$exp(j2πfot). В цифровой форме частотное смещение выражается в виде нормированного смещения $\in= f_oT$, где T - символьный интервал,
% длительность символа.
% Распространение сигнала по радиоканалу как по линейной системе описывается выражением свертки. Принятый сигнал в дискретной форме ( отсчеты
% сигнала) на выходе согласованного фильтра в отсчетные моменты времени записываются в виде


% $$y[n] = e^{i2\pi f_ot}\sum_{l=0}^{L}h[l]s[n-l] + v[n]$$

% h[n] - отсчеты комплексной импульсной характеристики радиоканала. При
% передаче в условиях помещения длина импульсной характеристики обычно небольшая, L = 1, 2. Предполагается, что символьная синхронизация уже установлена.
% Если величину частотного смещения оценить (измерить, численно определить по каком либо алгоритму), то влияние частотного смещения в отсчетах
% сигнала на выходе согласованного фильтра можно компенсировать, выполнив
% коррекцию. Если $\dot{\in}$ - оценка частотного смещения, то коррекция выполняется в
% виде умножения

% $$\dot{y}[n] =  e^{i2\pi \dot{\in}t}*y[n]$$

% \section{Решение (фазовая синхронизация)}

% Для оценки частотного смещения можно использовать метод, основанный на
% измерении изменения фазы в отсчётах преамбулы. Пreamble (преамбула) —
% это часть кадра данных, состоящая из двух одинаковых наборов отсчётов,
% причём начало преамбулы точно известно на приёмной стороне. Длина
% преамбулы составляет $N_t$ отсчётов.

% Сигнал преамбулы формируется в виде
% \begin{equation}
% s[n] = s[n + N_t] = t[n], \quad 0 \le n \le N_t - 1,
% \end{equation}
% где $t[n]$ — символы преамбулы.

% После приёма сигнала и выполнения символьной синхронизации отсчёты
% сигнала на выходе согласованного фильтра (СФ) имеют вид
% \begin{equation}
% y[n] = e^{j 2 \pi \varepsilon n} \sum_{l=0}^{L} h[l]\, s[n-l] + v[n],
% \end{equation}
% \begin{equation}
% \begin{aligned}
% y[n + N_t] &= e^{j 2 \pi \varepsilon (n + N_t)}
% \sum_{l=0}^{L} h[l]\, s[n + N_t - l] + v[n + N_t] \\
% &= e^{j 2 \pi \varepsilon N_t} \, y[n].
% \end{aligned}
% \end{equation}

% Таким образом, отсчёты сигнала преамбулы на выходе согласованного
% фильтра, разделённые интервалом $N_t$ отсчётов, отличаются на величину
% фазового сдвига
% \begin{equation}
% e^{j 2 \pi \varepsilon N_t},
% \end{equation}
% вызванного наличием частотного смещения.

% Оценка (измеренное значение) нормированного частотного смещения
% $\hat{\varepsilon}$ вычисляется по формуле
% \begin{equation}
% \hat{\varepsilon} =
% \frac{1}{2 \pi N_t}
% \arg\!\left(
% \sum_{l=L}^{N_t-1} y[l + N_t] \, y^{*}[l]
% \right),
% \end{equation}
% где $\arg(\cdot)$ — аргумент комплексного числа, а $y^{*}[n]$ —
% комплексно-сопряжённое значение сигнала на выходе согласованного
% фильтра. Рекомендуется исследовать влияние параметра $L$, принимая
% значения $L = 1, 2$.

% После вычисления оценки $\hat{\varepsilon}$ её необходимо применить к
% отсчётам сигнала данных для коррекции частотного смещения, после чего
% следует повторно построить временную диаграмму сигнала и диаграмму
% созвездия.


\endinput
